\documentclass[11pt]{article}

% Prefix for numedquestion's
\newcommand{\questiontype}{Question}


% Use this if your "written" questions are all under one section
% For example, if the homework handout has Section 5: Written Questions
% and all questions are 5.1, 5.2, 5.3, etc. set this to 5
% Use for 0 no prefix. Redefine as needed per-question.
\newcommand{\writtensection}{0}

\usepackage{amsmath, amsfonts, amsthm, amssymb}  % Some math symbols
\usepackage{mathtools}

\usepackage{centernot}
\usepackage{mathtools}

\usepackage{enumitem}

\setlength{\parindent}{0pt}

\begin{document}

\section{Goldwasser-Micali Cryptosystem}
A \textbf{Group Homomorphic Group Operation} is a group operation over group $\mathbb{G}$ that follows the following property:
$$h(x) \circ h(y) = h(x \circ y)$$

\subsection{Setup}
This cryptosystem is performed over a finite field $F_N$ where $N = p \cdot q$ for two large prime numbers $p,q$. $p,q$ are the secret keys and the $pk$ is $N$.

\subsection{Quadratic Residues}
An element $x$ is a \textbf{Quadratic Residue} mod $N$ if $\exists y \in Z$ such that $y^2 \mod N \equiv x$. The following facts follow:
\vspace{1em}

\textbf{Fact 1:} $\forall x$, $x$ is a quadratic residue mod $N$ iff $x$ is a quadratic residue mod $p$ AND $x$ is a quadratic residue mod $q$.
\bigskip

\textbf{Fact 2:} Let $p$ be an odd prime. They, the value of JACOBI symbol $J_p(x)$ for all $x$ is defined as:
$$J_p(x) = (x^{\frac{p-1}{2}} \mod p)$$
This is not obvious, but for any number in the group, the only possible values for the Jacobi Symbol is either $1$ or $-1$ for prime order groups.
\bigskip

\textbf{Fact 3:} Let $J_p(x)$ denote the value of Jacobi symbol on $x$ with an odd prime $p$, Then, $x$ is a quadratic residue $\mod p$ iff $J_p(x) = 1$.
\bigskip

\textbf{Fact 4:} If $N$ is a composite of two odd prime numbers $p$ and $q$ then
$$J_N(x) = J_p(x) \cdot J_q(x)$$
\bigskip

The assumption about Quadratic Residuity is that given $N=p \cdot q$ for sufficiently large primes $p$ and $q$, no PPT algorithm can efficiently determine whether $x$ is a quadratic residue mod $N$.

\subsection{Protocol}
The protocol is relatively simple. To encrypt an $m$, output a randomly chosen quadratic residue mod $N$ for $b=0$ and a non-residue for $b=1$. For decryption, the secret keys are used to compute $J_q(c)$ and $J_p(c)$ to determine if the cipher $c$ is a $1$ or $0$.
\bigskip

There are a bunch more notes about the number theory / math behind the cryptosystem but the main ideas were covered so I am skipping them here.
\bigskip

The punchline is the above cryptosystem has the following property an the ciphertexts of an encryptions $E(y),E(x)$.
$$E(x \oplus y) = E(x) \cdot E(y)$$
This means that if we multiply two group elements that are either QR or QNR, then the resulting element will represent the XOR of the two encrypted elements. More formally:
\begin{itemize}
    \item A QR $\cdot$ QNR is a QNR
    \item A QNR $\cdot$ QNR is a QR
    \item A QR $\cdot$ QR is a QR
\end{itemize}

\section{Private Information Retrieval}
Imagine a database $D$ that is $N$ bits long for some very large value $N$. Now imagine a user $u$ who wants to retrieve the ith bit $x_i$ from the database. However, $u$ does not want the database to learn what bit it requested. A trivial approach to this would be sending the entire database to the user, however, this is very expensive as it requires to send $N$ bits over some channel. The following is a protocol that uses $O(\sqrt{n})$ communication complexity that achieves this retrieval.

\subsection{Protocol}
This protocol requires an encryption scheme that that is fully homorphic that follows the following homomorphism:
$$E(x) \cdot E(y) = E(x \oplus y)$$
The protocol is as follows:
\begin{enumerate}
    \item The database $D$ will split the entire database into $\sqrt(N)$ length blocks so there are $\sqrt{n}$ blocks of length $\sqrt{n}$
    \item For each block, $U$ sends two values representing either a $0$ or a $1$ as $E(0)_j, E(1)_j$, encryptions of both values. However, for all blocks that do not contain the index of the bit $U$ wants, it actually just sends two encryptions of $0$.
    \item For every bit in the DB, $D$ replaces each bit with $E(b_i)$ sent by the user. Observe that for all bits except for the $\sqrt{N}$ that exist in the block that contains the $i$th bit, $D$ is replacing each bit with an encryption of $0$.
    \item Now, $D$ calculates, for the $j$th bit in each block, $$E(b_j) \oplus E(b_{j+\sqrt{n}})... \oplus E(b_{j+n-\sqrt{n}})$$ That is, the group operator on all encrypted bits at the $j$th index of each block. Since all blocks only encrypt $0$ except for the one with $i$, the above protocol will generate that block exactly since it becomes encryption:
    $$E(0 + 0 ... + E(b_{j+k\sqrt{n}})... + 0)$$
    \item $D$ sends this back to $U$, who checks the $i$th index to get its desired bit.
\end{enumerate}

\subsection{$PIR \Longrightarrow $ Collision Resistant Hashing}
The following is a proof that any non-trivial PIR, that is one that has sublinear communication complexity, implies a collision resistant hash. The protocol is as follows:
\bigskip

\textbf{Protocol}: Fix some index $i$ as the PIR query for database $x$ as well as the query for $i$. Then $h(x)$ is simply the PIR response.
\bigskip

Suppose that this scheme is not secure, meaning that an adversary can find two inputs $x,y$ such that $h(x) = h(y)$. This means that the $i$th bit must be the same in both $x$ and $y$ and that there is some index $j$ that is different. This immediately gives the adversary some information about the requested index $i$ meaning that the PIR is insecure, a contrdiction.

\section{Multiparty Computation}
Consider two parties, Alice and Bob, who have inputs $x,y$ respectively, that wish to compute some function $f$ without revealing any information beyond $f(x,y)$. It can be accomplished for the following gates as follows:
\bigskip

\textbf{XOR}: Reveal $x$ and $y$ to get $x \oplus y$. This is okay since $f(x,y)$ immediately reveals the inputs to both parties so no extra information is leaked.
\bigskip

\textbf{AND:} This cannot be done without some sort of cryptographic protocol, so we define the notion of a $1-2$ Oblivious Transfer:
\subsection{1-2 OT}
This device defined by [Even, Goldreich, Lempel] has two parties, a sender $S$ and receiver $R$ and acts as follows:
\begin{enumerate}
    \item $S$ provides two bits as input $b_0$ and $b_1$ and inputs it into the $1-2$ OT.
    \item $R$ can query the OT one time with $s$ that is either $0$ or $1$ and receives $b_s$.
\end{enumerate}
Now, we can design a secure protocol for computing AND as follows:
\begin{enumerate}
    \item Alice inputs $b_0 = x \land 0$ and $b_1 = x \land 1$.
    \item Bob inputs $s$ as $y$ to get $y \land x$ and outputs it to Alice.
\end{enumerate}

\subsection{Multi Party XOR/AND computation}
Consider $n$ parties who wish to compute the $\oplus$ of their bits $b_1,..., b_n$. This can be accomplished if the parties can communicate pairwise with one another. It can be accomplished in the following manner.
\begin{enumerate}
    \item Player 1 $p_1$ xors their bits with a completely random bitstring to get $m_1 = x_1 \oplus r$.
    \item Each subsequent player calculates their message $m_i = x_{i} \oplus m_{i-1}$ for $i=2...n$.
    \item Player $n$ sends $m_n$ to $p_1$ who calculates $r \oplus m_n = \bigoplus_{i=1}^{n} x_i$ and outputs the result to all other parties.
\end{enumerate}

\section{Secure Computation}
We discuss a secure $2$-Party Computation that can compute a circuit consisting of only XOR and AND gates. Two do this, there is a notion of secret sharing that inductively implies throughout the circuit from input to output. We also assume that the circuit outputs either $0/1$, meaning it has one output wire at the top level.
\subsection{Secret Sharing}

Secret sharing involves splitting a secret value $S$ into shares $s_1,...,s_n$ distributed among $n$ players, given a certain number of shares, say $t$, one can recover the secret $S$.

\subsection{Setup}
We consider the two party case where parties Alice $A$ and Bob $B$ have inputs $X=x_1,...,x_n$ and $Y=y_1,...,y_n$ respsectively. They need to split their inputs into shares such that using those shares, they can compute the corresponding outputs to XOR and AND gates produce corresponding shares to the output of the gate. We split up $A$'s input into shares by selecting a random $a_i$ for each bit $x_i$ and creating share $a_i' = x_i \oplus a_i$. Now $A$'s share is $a_i$ and $B$'s share of $x_i$ is $a_i'$. This is also done for $Y$ to create shares $b_i$ and $b_i'$. Note that $A$ has shares that we will call $a_1, b_1$ and $B$ has shares $a_2, b_2$ and that $a_1 \oplus a_2 = x$ and $b_1 \oplus b_2 = y$.

\subsection{XOR Gates}
To simulate an XOR gate, we can simply create shares for $A$ and $B$ as:
$$a_s = a_1 \oplus b_1$$
$$b_s = a_2 \oplus b_2$$
This works since we are trying to compute the value of $a \oplus b$ which can be represented as $a_1 \oplus a_2 \oplus b_1 \oplus b_2$. This value is split into two shares above and therefore XOR gates can be simulated.
\subsection{AND gates}
To compute and gates, we need to compute $(a_1 \oplus a_2)(b_1 \oplus b_2)$. To accomplish this, we distribute to get:
$$(a_1 \oplus a_2)(b_1 \oplus b_2) = a_1b_1 \oplus a_1b_2 \oplus a_2b_1 \oplus a_2b_2$$
Recall that $A$ has shares $a_1,b_1$ and $B$ has shares $a_2,b_2$ so the first and last terms can be calculated. They calculate terms $a_1b_2$ and $a_2b_1$ in the following manner:
\begin{enumerate}
    \item $A$ uses 1-2 OT to send $s_0 = a_10 \oplus r$ and $s_1 = a_11 \oplus r$ to $B$ who inputs $s=b_2$ to obtain $a_1b_2 \oplus r$
    \item The same is done with $B$ to compute $a_2b_1 \oplus r'$ which is given to $A$
    \item $A$ generates share $r \oplus a_1b_1 \oplus a_2b_1 \oplus r'$
    \item $B$ generates share $r' \oplus a_2b_2 \oplus a_1b_2 \oplus r$
\end{enumerate}
The above shares clearly result in shares for $AB$ demonstrating the correctness of the AND gate.


\end{document}